%% ============================================================
%%  Predictive Risk Intelligence for Metabolic Screening in Diabetes — Overleaf-Ready LaTeX
%%  Compiler: pdfLaTeX  |  Engine: TeX Live 2023+
%% ============================================================
\documentclass[12pt,a4paper,oneside]{report}

%% ----- packages -----
\usepackage[T1]{fontenc}
\usepackage[utf8]{inputenc}
\usepackage{lmodern}
\usepackage[top=2.54cm, bottom=2.54cm, left=3.5cm, right=2.54cm]{geometry}
\usepackage{setspace}
\usepackage{graphicx}
\usepackage{booktabs}
\usepackage{longtable}
\usepackage{array}
\usepackage{amsmath,amssymb,amsfonts}
\usepackage{hyperref}
\usepackage{xcolor}
\usepackage{fancyhdr}
\usepackage{titlesec}
\usepackage{tocloft}
\usepackage{enumitem}
\usepackage{microtype}
\usepackage{caption}
\usepackage{listings}
\usepackage{float}
\usepackage{multirow}
\usepackage{tabularx}
\usepackage[most]{tcolorbox}  % equation boxes

%% ----- spacing (no parskip package, manual control) -----
\setlength{\parskip}{6pt plus 1pt minus 1pt}
\setlength{\parindent}{0pt}

%% ----- prevent text cut at page bottom -----
\raggedbottom
\widowpenalty=10000
\clubpenalty=10000
\displaywidowpenalty=10000

%% ----- hyperref setup -----
\hypersetup{
    colorlinks=true,
    linkcolor=black,
    citecolor=black,
    urlcolor=blue,
    pdftitle={Predictive Risk Intelligence for Metabolic Screening in Diabetes},
    pdfauthor={},
    pdfkeywords={Identifiable VAE, Conformal Prediction, Metabolic Phenotyping, Normal-BMI MASLD, Semi-Supervised Learning}
}

%% ----- line spacing -----
\doublespacing

%% ----- header/footer -----
\pagestyle{fancy}
\fancyhf{}
\fancyhead[L]{\small\leftmark}
\fancyhead[R]{\small\thepage}
\fancyfoot{}
\renewcommand{\headrulewidth}{0.4pt}
%% fix chapter-opening pages so they show page number, not nothing
\fancypagestyle{plain}{%
  \fancyhf{}
  \fancyfoot[C]{\thepage}
  \renewcommand{\headrulewidth}{0pt}%
}

%% ----- chapter title style -----
\titleformat{\chapter}[display]
  {\normalfont\huge\bfseries}{\chaptername\ \thechapter}{20pt}{\Huge}
\titlespacing*{\chapter}{0pt}{-10pt}{40pt}

%% ----- listings style (for code) -----
\lstset{
  basicstyle=\ttfamily\small,
  breaklines=true,
  frame=single,
  keywordstyle=\bfseries,
  commentstyle=\itshape\color{gray},
}

%% ----- equation highlight box -----
\tcbuselibrary{skins,breakable}
\newtcolorbox{equationbox}{
  enhanced,
  breakable,
  colback=blue!4!white,
  colframe=blue!45!black,
  boxrule=0pt,
  leftrule=3pt,
  arc=2pt,
  left=10pt, right=10pt, top=6pt, bottom=6pt,
  before skip=10pt,
  after skip=10pt,
}

%% ----- convenience commands -----
\newcommand{\lmsis}{\textsc{lmsis}}
\newcommand{\daivae}{DA-SS-iVAE}
\newcommand{\rhoval}[1]{\rho = #1}
\newcommand{\auroc}[1]{\textsc{auroc} = #1}

%% ============================================================
\begin{document}
\pagenumbering{roman}

%% ============================================================
%%  TITLE PAGE
%% ============================================================
\begin{titlepage}
\centering
\vspace*{1cm}

{\Large\bfseries
PREDICTIVE RISK INTELLIGENCE FOR METABOLIC SCREENING IN DIABETES\\[6pt]
{\large A Semi-Supervised Deep Learning Framework for Early Detection of\\[2pt]
Metabolic Dysfunction in Normal-BMI Adults}
\par}

\vspace{1.5cm}
\textit{A Dissertation submitted in partial fulfilment of the requirements\\
for the degree of Master of Computer Applications}

\vspace{1.5cm}
\textbf{Submitted by}\\[8pt]
\begin{tabular}{ll}
\underline{\hspace{5.5cm}} & \underline{\hspace{5.5cm}} \\
\textit{(Name \& Registration Number)} & \textit{(Name \& Registration Number)} \\[10pt]
\underline{\hspace{5.5cm}} & \underline{\hspace{5.5cm}} \\
\textit{(Name \& Registration Number)} & \textit{(Name \& Registration Number)}
\end{tabular}

\vspace{1.5cm}
\textbf{Under the Supervision of}\\[6pt]
\underline{\hspace{8cm}}\\[4pt]
\textit{(Supervisor Name \& Designation)}

\vspace{1.5cm}
\textbf{Department of Computer Applications}\\
\underline{\hspace{8cm}}\\[6pt]
\textbf{June 2026}

\end{titlepage}

%% ============================================================
%%  CERTIFICATE PAGE
%% ============================================================
\clearpage
\chapter*{Certificate}
\thispagestyle{plain}
\addcontentsline{toc}{chapter}{Certificate}

This is to certify that the dissertation entitled

\begin{quote}
\textbf{``Predictive Risk Intelligence for Metabolic Screening in Diabetes''}
\end{quote}

\noindent submitted by \underline{\hspace{5cm}} (Registration No.\ \underline{\hspace{3cm}}) in partial fulfilment of the requirements for the degree of Master of Computer Applications has been carried out under my supervision. The work reported in this dissertation is original and has not been submitted elsewhere for any degree.

\vspace{2cm}
\begin{flushright}
\rule{7cm}{0.4pt}\\
\underline{\hspace{6cm}}\\
\textit{(Supervisor Name \& Designation)}\\
Department of Computer Applications\\
\underline{\hspace{5cm}}\\
Date: \underline{\hspace{3cm}}
\end{flushright}

%% ============================================================
%%  DECLARATION
%% ============================================================
\clearpage
\chapter*{Declaration}
\addcontentsline{toc}{chapter}{Declaration}

We hereby declare that the work presented in this dissertation, titled \textbf{``Predictive Risk Intelligence for Metabolic Screening in Diabetes''}, submitted to \underline{\hspace{5cm}} for the award of the degree of Master of Computer Applications, is our collective original work. The dissertation has not been submitted for any other degree or examination in this or any other university.

All sources of information used in this work have been duly cited and referenced. Any ideas or data taken from other sources have been acknowledged according to standard referencing practices.

\vspace{1.5cm}
\begin{flushleft}
\begin{tabular}{p{7cm} p{7cm}}
\underline{\hspace{6cm}} & \underline{\hspace{6cm}} \\
\textit{(Name \& Registration Number)} & \textit{(Name \& Registration Number)} \\[20pt]
\underline{\hspace{6cm}} & \underline{\hspace{6cm}} \\
\textit{(Name \& Registration Number)} & \textit{(Name \& Registration Number)} \\[12pt]
\end{tabular}\\
Date: \underline{\hspace{3cm}}
\end{flushleft}

%% ============================================================
%%  ACKNOWLEDGEMENTS
%% ============================================================
\clearpage
\chapter*{Acknowledgements}
\addcontentsline{toc}{chapter}{Acknowledgements}

We would like to express our sincere gratitude to our supervisor for their guidance, patience, and constructive feedback throughout this project. Their advice helped us approach the clinical and machine learning aspects of this problem with honesty and rigour.

We are grateful to our department and university for providing the computational resources and academic environment necessary to pursue this research. We also thank our families and friends for their constant support.

This work makes use of publicly available data from the National Health and Nutrition Examination Survey (NHANES), provided by the Centers for Disease Control and Prevention (CDC). We thank the survey participants whose anonymised data made this research possible.

%% ============================================================
%%  ABSTRACT
%% ============================================================
\clearpage
\chapter*{Abstract}
\addcontentsline{toc}{chapter}{Abstract}

Body Mass Index (BMI) is the predominant screening metric in primary care cardiometabolic assessment. However, in adults with normal BMI (18.5--24.9\,kg/m\textsuperscript{2}), this proxy offers limited discriminative utility regarding fat distribution, visceral adiposity, and ectopic lipid accumulation. Some normal-weight individuals may exhibit concurrent insulin resistance and hepatic steatosis without triggering standard screening thresholds---a combination associated with increased cardiovascular and hepatic disease risk.

This dissertation, Predictive Risk Intelligence for Metabolic Screening in Diabetes, presents the Latent Metabolic State Inference System (LMSIS), a semi-supervised deep learning framework for early detection of metabolic dysfunction in normal-BMI adults. The core model, the \textbf{Dual-Anchored Semi-Supervised Identifiable Variational Autoencoder (\daivae{})}, is designed according to Nonlinear ICA identifiability principles (up to coordinate permutation and element-wise monotone scaling) through demographic-conditioned priors and regularises biological interpretability through monotone anchor networks aligned to reference-standard measurements: HOMA-IR for insulin resistance ($Z_1$) and FibroScan Controlled Attenuation Parameter (CAP) for hepatic steatosis ($Z_2$).

Trained exclusively on the NHANES J-cycle (2017--2018, $n_{\text{train}} = 401$, random 70\% of J-cycle complete records), the DA-SS-iVAE was evaluated on the strict temporal holdout P-cycle (2019--2020), which contains zero training participants. A five-seed robustness analysis yields a mean P-cycle Spearman correlation of $\rho = 0.583 \pm 0.027$ against FibroScan CAP ($n = 870$; all five seeds exceed the mixed-split baseline of 0.501). An AUROC of $0.841$ is achieved when the latent steatosis coordinate $Z_2$ is evaluated against FibroScan CAP thresholds (CAP $\geq 248$\,dB/m). Mondrian conformal calibration restores $\geq 90\%$ subgroup coverage for the highest-risk Dual-Burden quadrant, which achieves only 81.6\% coverage under standard marginal calibration.

All results are reported with their full uncertainty ranges, sample-size limitations, and simulation-induced inflation disclaimers. The system is implemented as a full-stack research prototype comprising a PyTorch inference engine, a FastAPI backend service, and an interactive React-based Research Translation Layer.

\vspace{0.5cm}
\noindent\textbf{Keywords:} Identifiable Variational Autoencoder, Conformal Prediction, Metabolic Phenotyping, Normal-BMI MASLD, Semi-Supervised Learning, Symbolic Regression, Riemannian Geometry, Clinical Machine Learning Systems.

%% ============================================================
%%  TABLE OF CONTENTS / LIST OF TABLES / LIST OF FIGURES
%% ============================================================
\clearpage
\tableofcontents

\clearpage
\listoftables
\addcontentsline{toc}{chapter}{List of Tables}

\clearpage
%% (add \listoffigures if you include figures)

%% ============================================================
%%  LIST OF ABBREVIATIONS
%% ============================================================
\clearpage
\chapter*{List of Abbreviations}
\addcontentsline{toc}{chapter}{List of Abbreviations}
\begin{tabular}{@{}ll@{}}
\toprule
\textbf{Abbreviation} & \textbf{Expansion} \\
\midrule
AIP       & Atherogenic Index of Plasma \\
ALT       & Alanine Aminotransferase \\
AST       & Aspartate Aminotransferase \\
AUROC     & Area Under the Receiver Operating Characteristic Curve \\
BMI       & Body Mass Index \\
BUN       & Blood Urea Nitrogen \\
CAP       & Controlled Attenuation Parameter \\
CDC       & Centers for Disease Control and Prevention \\
\daivae{} & Dual-Anchored Semi-Supervised Identifiable Variational Autoencoder \\
DCR       & Discriminative Contribution Ratio \\
ELBO      & Evidence Lower Bound \\
FLI       & Fatty Liver Index \\
GGT       & Gamma-Glutamyl Transferase \\
HDL       & High-Density Lipoprotein \\
HOMA-IR   & Homeostatic Model Assessment for Insulin Resistance \\
HSI       & Hepatic Steatosis Index \\
ICA       & Independent Component Analysis \\
IR        & Insulin Resistance \\
iVAE      & Identifiable Variational Autoencoder \\
LDL       & Low-Density Lipoprotein \\
\lmsis{}  & Latent Metabolic State Inference System \\
MASLD     & Metabolic Dysfunction-Associated Steatotic Liver Disease \\
MASH      & Metabolic Dysfunction-Associated Steatohepatitis \\
MetS      & Metabolic Syndrome \\
MHNW      & Metabolically Healthy Normal Weight \\
MLP       & Multi-Layer Perceptron \\
NAFLD     & Non-Alcoholic Fatty Liver Disease \\
NHANES    & National Health and Nutrition Examination Survey \\
OOD       & Out-of-Distribution \\
PySR      & Python Symbolic Regression \\
TyG       & Triglyceride-Glucose Index \\
VAE       & Variational Autoencoder \\
VCTE      & Vibration-Controlled Transient Elastography \\
\bottomrule
\end{tabular}

%% ============================================================
%%  MAIN MATTER
%% ============================================================
\clearpage
\pagenumbering{arabic}
\setcounter{page}{1}

%% ============================================================
%%  CHAPTER 1: INTRODUCTION
%% ============================================================
\chapter{Introduction}

\section{Introduction}

Body Mass Index (BMI) is the predominant screening metric in primary care cardiometabolic assessment. While BMI correlates with total adiposity in obese populations, it offers limited discriminative utility regarding regional fat distribution, visceral adiposity, and ectopic lipid accumulation. In adults with normal BMI (18.5--24.9\,kg/m\textsuperscript{2}), this proxy may fail to identify biochemically hidden metabolic risk factors: a substantial proportion may exhibit concurrent insulin resistance and hepatic steatosis---conditions associated with increased cardiovascular and hepatic disease risk in the absence of standard screening alerts.

This dissertation, Predictive Risk Intelligence for Metabolic Screening in Diabetes, presents the Latent Metabolic State Inference System (LMSIS), a semi-supervised deep learning framework for early detection of metabolic dysfunction in normal-BMI adults. The core model---the Dual-Anchored Semi-Supervised Identifiable Variational Autoencoder (\daivae{})---is designed according to Nonlinear ICA identifiability principles (up to coordinate permutation and element-wise monotone scaling) through demographic-conditioned priors and regularises biological interpretability through monotone anchor networks aligned to reference-standard measurements: HOMA-IR for insulin resistance ($Z_1$) and FibroScan Controlled Attenuation Parameter (CAP) for hepatic steatosis ($Z_2$).

The system is implemented as a full-stack research prototype comprising a PyTorch-based inference engine, a FastAPI backend service, and an interactive React-based Research Translation Layer. All components are version-controlled, integration-tested, and publicly documented. This work is situated at the intersection of deep generative modelling, conformal uncertainty quantification, and clinical machine-learning systems engineering.

\section{Existing Systems}

Contemporary metabolic screening relies on two categories of tools: clinical scoring indices and supervised machine-learning classifiers.

\paragraph{Clinical Scoring Indices.} The Fatty Liver Index (FLI), Hepatic Steatosis Index (HSI), Triglyceride-Glucose Index (TyG), and NAFLD Liver Fat Score (LFS) are widely used in primary care. These indices combine BMI, liver enzymes, lipid panels, and glucose metrics into scalar risk scores. They are inexpensive, interpretable, and validated in large mixed-BMI cohorts.

\paragraph{Machine Learning Classifiers.} Recent work applies deep neural networks and gradient-boosted trees to wearable sensor data \citep{metwally2026} and NHANES biomarkers \citep{zhang2025,tertulino2025} for MASLD risk estimation. Unsupervised approaches such as the SENA-$\delta$ VAE \citep{shen2025} model causal discrepancies in genetic perturbation data but do not anchor latent factors to clinically measurable phenotypes.

Among the systems identified in the reviewed literature, all share a critical limitation when restricted to normal-BMI populations: they are calibrated on mixed-BMI cohorts where BMI contributes substantial discriminative variance. When applied to normal-BMI adults, the reduced variance of BMI diminishes the predictive signal of these models, producing degraded or inverted risk rankings.

\section{Problem Statement}

Standard cardiometabolic screening workflows use BMI as a primary gatekeeper. In normal-BMI adults, this gate may fail to identify a population with genuine metabolic risk. Existing clinical indices and ML classifiers, trained on mixed-BMI distributions, experience predictive degradation when BMI variance is restricted. The specific problems are:

\begin{enumerate}[leftmargin=*]
    \item \textbf{Discriminative Signal Loss:} Clinical indices such as HSI depend linearly on BMI; when BMI is restricted to a narrow band, their discriminative contribution ratio is substantially reduced, rendering the index less informative.

    \item \textbf{Risk Ranking Inversion:} Indices such as NAFLD-LFS incorporate metabolic syndrome criteria that correlate positively with liver fat in obese populations but not necessarily in normal-BMI adults, where steatosis may occur via distinct non-obese pathways (ectopic lipid deposition, insulin-mediated de novo lipogenesis). This can produce ranking inversion, where higher-risk patients receive lower scores.

    \item \textbf{Opacity of Deep Models:} Standard variational autoencoders suffer from latent space non-identifiability---any rotation of the latent axes yields identical reconstruction loss, leaving axes biologically uninterpretable.

    \item \textbf{Subgroup Coverage Failure:} Standard (marginal) conformal prediction provides only global coverage properties. High-risk subgroups with shifted covariate distributions may experience systematic under-coverage, leaving vulnerable patients insufficiently protected by uncertainty bounds.
\end{enumerate}

\paragraph{Motivation and Significance.} From a computer science perspective, this problem domain offers a rigorous testbed for identifiable generative modelling, semi-supervised learning under structured missingness, and subgroup-safe uncertainty quantification. From a clinical research perspective, a robust, interpretable, and uncertainty-aware computational inference prototype for normal-BMI metabolic risk mapping serves to demonstrate primary care screening translation possibilities.

\section{Research Gap}

Based on the analysis of existing systems, four specific research gaps are identified:

\begin{enumerate}[leftmargin=*]
    \item \textbf{Identifiable Latent Spaces in Clinical ML.} Existing VAE-based clinical models do not enforce formal identifiability. There is a gap for a deep generative model that recovers biologically meaningful latent coordinates up to permutation and monotone transformation under theoretical conditions \citep{khemakhem2020,kingma2014}.

    \item \textbf{Normal-BMI Specific Modelling.} Among the systems identified in the reviewed literature, all contemporary clinical ML systems for MASLD screening are trained and validated on mixed-BMI cohorts. To the best of the authors' knowledge following the reviewed literature, no identified system has been specifically designed, trained, and validated exclusively for normal-BMI adults.

    \item \textbf{Subgroup-Safe Uncertainty Quantification.} Existing clinical ML models report point predictions or global confidence intervals. Among the systems reviewed, none provide subgroup-conditional coverage properties \citep{barber2023,vovk2005}.

    \item \textbf{Interpretable Deep Generative Models with Closed-Form Extraction.} There is a gap for methods that extract human-interpretable closed-form relationships from learned neural decoders without sacrificing representational capacity.
\end{enumerate}

\section{Objectives}

The primary objective of this dissertation is to design, implement, and validate a clinical machine-learning pipeline that infers concurrent metabolic risk from routine blood biomarkers in normal-BMI adults. The specific objectives are:

\begin{enumerate}[leftmargin=*]
    \item To design and implement a \daivae{} according to theoretical identifiability conditions (up to coordinate permutation and scaling) and biological interpretability through demographic-conditioned prior networks and monotone anchor networks.
    \item To implement Mondrian conformal prediction stratified by phenotypic quadrant, restoring subgroup-conditional coverage guarantees for high-risk subpopulations.
    \item To extract interpretable closed-form equations from the learned decoder using symbolic regression, enabling hypothesis generation about biomarker-latent relationships.
    \item To develop a computational method for optimal path planning on the learned Riemannian manifold.
    \item To engineer a full-stack research prototype (FastAPI backend service, React-based Research Translation Layer) with integration testing, API documentation, and reproducible development pipelines.
\end{enumerate}

\section{Methodology and Scope}

\paragraph{Methodology.} This study employs a pipeline methodology comprising:
\begin{itemize}[leftmargin=*]
    \item \textbf{Data Source:} NHANES 2017--2020, providing nationally representative biomarker, demographic, and FibroScan VCTE data.
    \item \textbf{Model Architecture:} A custom PyTorch implementation of the \daivae{} with residual encoder, conditional prior networks, monotone anchor networks, and residual decoder.
    \item \textbf{Training Regime:} Semi-supervised learning with binary masking for missing CAP labels, ELBO maximisation with auxiliary anchor losses and Frobenius covariance penalty.
    \item \textbf{Uncertainty Quantification:} Split conformal prediction and Mondrian (stratified) conformal prediction across four phenotypic quadrants.
    \item \textbf{Interpretability:} Symbolic regression via PySR to extract closed-form decoder approximations.
    \item \textbf{Computational Geometry:} Riemannian metric tensor computation via PyTorch Autograd, with Euler-Lagrange geodesic solving using Brent's root-finding method.
    \item \textbf{System Engineering:} FastAPI backend (Python 3.12), React 19 Research Translation Layer (Vite, Tailwind CSS, D3.js, Framer Motion), pytest integration suite.
\end{itemize}

\paragraph{Scope.} The study is restricted to:
\begin{itemize}[leftmargin=*]
    \item Normal-BMI adults ($18.5 \leq \text{BMI} \leq 24.9$\,kg/m\textsuperscript{2}) aged 18+ with complete biomarker panels.
    \item Cross-sectional inference only; longitudinal trajectory modelling is excluded.
    \item A research prototype; clinical deployment, regulatory clearance (FDA/CE), and prospective outcome validation are explicitly out of scope.
    \item 14 routine blood biomarkers as input features; imaging or genomic data are not incorporated.
\end{itemize}

\section{Summary of Research Contributions}

This dissertation presents the following core research contributions:

\begin{enumerate}[leftmargin=*]
    \item \textbf{Normal-BMI Specific Metabolic Modelling Pipeline:} A custom screening and metabolic representation pipeline specifically for normal-weight adults ($18.5 \leq \text{BMI} \leq 24.9$\,kg/m\textsuperscript{2}), addressing the discriminative collapse and risk-inversion failures of standard clinical indices when applied to normal-BMI cohorts.
    \item \textbf{\daivae{} Architecture:} A model designed according to the theoretical iVAE identifiability conditions of Khemakhem et al.\ (2020), using demographic-conditioned priors and positive-weight monotone anchor networks aligned with HOMA-IR and CAP, with the aim of recovering approximately identifiable latent coordinates.
    \item \textbf{Mondrian Subgroup Calibration:} Mondrian conformal prediction stratified by metabolic quadrant, restoring risk interval coverage to $\geq 90\%$ for the high-risk Dual-Burden subpopulation, which suffers from systematic coverage collapse (81.6\%) under standard global marginal calibration.
    \item \textbf{Symbolic Regression Interpretability Layer:} PySR applied to extract human-readable, closed-form algebraic approximations from the deep residual decoder.
    \item \textbf{Research Prototype Implementation:} A complete, reproducible research system including FastAPI inference endpoints, a Riemannian metric tensor solver, and an interactive React-based Research Translation Layer.
\end{enumerate}

\section{Chapter Summary}

This chapter introduced the clinical and technical motivation for LMSIS, identified four specific research gaps in existing clinical ML systems, and defined five objectives aligned with identifiable generative modelling, subgroup-safe uncertainty quantification, interpretability extraction, and full-stack systems engineering. The following chapters review the technical background, compare existing methods, present the proposed \daivae{} architecture and system implementation, and conclude with limitations and future extensions.

%% ============================================================
%%  CHAPTER 2: BACKGROUND AND LITERATURE REVIEW
%% ============================================================
\chapter{Background and Literature Review}

\section{Domain Overview and Definitions}

\paragraph{Metabolic Dysfunction-Associated Steatotic Liver Disease (MASLD).} Formerly known as NAFLD, MASLD is defined as hepatic steatosis ($\geq 5\%$ fat accumulation) in the absence of excessive alcohol consumption, coupled with metabolic risk factors. The disease spectrum ranges from simple steatosis to steatohepatitis (MASH), fibrosis, and cirrhosis.

\paragraph{Insulin Resistance (IR).} A pathological condition in which cells fail to respond normally to insulin. In clinical practice, IR is commonly approximated by the Homeostatic Model Assessment for Insulin Resistance (HOMA-IR):
\begin{equationbox}
\[
\text{HOMA-IR} = \frac{\text{fasting glucose} \times \text{fasting insulin}}{22.5}
\]\vspace{-6pt}
\end{equationbox}

\paragraph{Hepatic Steatosis Quantification.} The Controlled Attenuation Parameter (CAP) from FibroScan VCTE (Vibration-Controlled Transient Elastography) provides a non-invasive, quantitative measure of liver fat accumulation, reported in decibels per metre (dB/m). CAP $\geq 248$\,dB/m indicates clinically significant steatosis.

\paragraph{The Normal-BMI Phenotype.} While MASLD is traditionally associated with obesity, approximately 10--20\% of cases occur in normal-BMI individuals (``lean MASLD''). In this population, steatosis is driven by ectopic lipid deposition, visceral adiposity not captured by BMI, and insulin-mediated de novo lipogenesis rather than by total body mass.

\paragraph{Identifiability in Latent Variable Models.} A latent variable model is identifiable if the true latent factors can be recovered up to permutation and element-wise invertible transformation. Standard VAEs are non-identifiable because any orthogonal transformation of the latent space preserves the ELBO. The iVAE framework \citep{khemakhem2020} proves identifiability under the condition that the prior is conditioned on sufficiently informative auxiliary variables.

\section{Traditional Clinical Approaches}

\paragraph{Body Mass Index (BMI).} $\text{BMI} = \text{weight (kg)} / \text{height}^2 \text{ (m}^2\text{)}$. The WHO defines normal BMI as 18.5--24.9\,kg/m\textsuperscript{2}. While BMI is a population-level correlate of adiposity, it does not distinguish muscle from fat, nor subcutaneous from visceral fat.

\paragraph{Clinical Steatosis Indices.}
\begin{itemize}[leftmargin=*]
    \item \textbf{Fatty Liver Index (FLI):} Combines BMI, waist circumference, triglycerides, and GGT. Validated in mixed-BMI Italian cohorts.
    \item \textbf{Hepatic Steatosis Index (HSI):} $\text{HSI} = 8 \times (\text{ALT}/\text{AST}) + \text{BMI} + \text{adjustment terms}$. Heavily dependent on BMI as a linear discriminant.
    \item \textbf{NAFLD Liver Fat Score (LFS):} Incorporates metabolic syndrome criteria, type 2 diabetes status, and liver enzymes.
    \item \textbf{TyG Index:} A simple insulin resistance proxy derived from fasting triglycerides and glucose.
\end{itemize}

These indices are inexpensive and interpretable but share a structural dependency on BMI that becomes problematic when BMI variance is restricted.

\section{Existing Machine Learning Approaches}

\paragraph{Supervised Classification.} \citet{zhang2025} and \citet{tertulino2025} apply XGBoost and Random Forest to NHANES biomarkers for binary MASLD classification. These models achieve moderate AUROC in mixed-BMI cohorts but are not validated in normal-BMI subgroups and produce scalar predictions without latent structure.

\paragraph{Wearable-Enhanced Deep Learning.} \citet{metwally2026} combine wearable sensor streams with deep neural networks to predict HOMA-IR in a mixed-BMI cohort ($n = 1{,}165$). The model does not address hepatic steatosis and provides no subgroup calibration guarantees.

\paragraph{Clustering-Based Subtyping.} \citet{lee2024} apply K-means and partitioning methods to liver biopsy histology in an obese-majority MASLD cohort. This work identifies discrete histological subtypes but does not provide continuous, anchored latent coordinates applicable to routine blood biomarkers.

\paragraph{Causal Discrepancy VAEs.} \citet{shen2025} propose the SENA-$\delta$ VAE for gene-perturbation response prediction. While methodologically sophisticated, it operates on in-vitro genetic assay data and does not anchor latent factors to clinical reference-standard measurements.

\begin{table}[H]
\centering
\caption{Positioning of LMSIS within contemporary clinical ML.}
\label{tab:positioning}
\resizebox{\textwidth}{!}{%
\begin{tabular}{@{}lllllll@{}}
\toprule
\textbf{Study / Method} & \textbf{Year} & \textbf{Core Algorithm} & \textbf{Population} & \textbf{Latent Topology} & \textbf{Validation} & \textbf{Subgroup Calibration} \\ \midrule
Metwally et al.\ (WEAR-ME) & 2026 & DNN + Wearables & Mixed-BMI & Scalar & Cross-validation & None \\
Zhang et al.\ (PLOS ONE) & 2025 & XGBoost/RF & Mixed-BMI & Binary & Train-Test & None \\
Lee et al.\ (Nat.\ Med.) & 2024 & K-means & Obese-majority & Discrete clusters & Liver biopsy & None \\
SENA-$\delta$ VAE & 2025 & Causal VAE & Genetic perturb-seq & Unanchored & In-vitro & None \\
\textbf{LMSIS (This work)} & \textbf{2026} & \textbf{\daivae{}} & \textbf{Normal-BMI specific} & \textbf{Continuous 2D anchored} & \textbf{OOD temporal + VCTE} & \textbf{Mondrian $\geq 90\%$} \\ \bottomrule
\end{tabular}%
}
\end{table}

\section{Literature Review Methodology}

This review employs a \textbf{narrative synthesis} approach rather than a formal systematic review (PRISMA/SLR), as the intersection of identifiable VAEs, normal-BMI metabolic phenotyping, and conformal prediction is too narrow for comprehensive PRISMA criteria. The search strategy targeted:
\begin{itemize}[leftmargin=*]
    \item Clinical ML for MASLD/NAFLD screening (2019--2026)
    \item Identifiable latent variable models (iVAE, nonlinear ICA)
    \item Conformal prediction in healthcare applications
    \item NHANES-based metabolic inference studies
\end{itemize}

\paragraph{Inclusion criteria:} Peer-reviewed publications or high-quality preprints with publicly described methodology, validation on human biomarker data, and relevance to metabolic phenotyping.

\paragraph{Exclusion criteria:} Purely genetic or animal studies, models without validation, and non-English sources without translation.

\section{Limitations of Existing Work}

The literature review reveals five consistent limitations across existing clinical ML systems:

\begin{enumerate}[leftmargin=*]
    \item \textbf{Mixed-BMI Calibration Bias.} All existing indices and classifiers are calibrated on cohorts where BMI spans the full range. Their performance in restricted normal-BMI subgroups is either unreported or degraded.
    \item \textbf{Unanchored Latent Spaces.} VAE-based approaches in clinical ML do not assign biological meaning to latent axes.
    \item \textbf{Absence of Subgroup Coverage Properties.} To the best of the authors' knowledge following the reviewed literature, no identified metabolic phenotyping system implements subgroup-conditional conformal calibration.
    \item \textbf{Opacity of Decoder Mappings.} Even when deep models are accurate, their internal biomarker-latent relationships remain hidden.
    \item \textbf{Narrow Technical Scope.} Existing works are either pure ML methods papers without system engineering, or pure software systems without methodological novelty.
\end{enumerate}

%% ============================================================
%%  CHAPTER 3: COMPARATIVE ANALYSIS
%% ============================================================
\chapter{Comparative Analysis}

\section{Comparative Framework}

This chapter provides a comparative analysis of existing metabolic screening methods and the proposed LMSIS pipeline across three dimensions:
\begin{itemize}[leftmargin=*]
    \item \textbf{Predictive Validity:} Correlation with reference-standard biomarkers (HOMA-IR, CAP) in normal-BMI cohorts.
    \item \textbf{Interpretability:} Whether the model provides biologically meaningful, axis-aligned latent coordinates.
    \item \textbf{Safety:} Whether the model provides subgroup-conditional uncertainty coverage.
\end{itemize}

All methods are evaluated on the same normal-BMI NHANES sub-cohort. The primary \daivae{} model is trained on a strict temporal split using J-cycle data only (2017--2018; $n_{\text{train}} = 401$ complete records representing 70\% of J-cycle, $n_{\text{val}} = 86$ validation, and $n_{\text{test}} = 87$ internal test), while the independent P-cycle cohort (2019--2020; $n = 903$) is held out entirely as a strict temporal out-of-distribution evaluation split. For comparative baseline analysis, evaluations are also reported on a baseline model trained using a random 70/30 split of the combined J+P cohort ($n = 1{,}033$ training complete records, $n = 444$ held-out validation/test). Classification AUROC is computed on the respective CAP-labelled subsets ($n = 1{,}422$ CAP-labelled total across J+P: $552$ in J-cycle, $870$ in P-cycle).

\section{Analysis of Clinical Scoring Indices in Normal-BMI Cohorts}

\subsection{HSI Discriminative Collapse}

The Hepatic Steatosis Index is defined as:
\[
\text{HSI} = 8 \times \left(\frac{\text{ALT}}{\text{AST}}\right) + \text{BMI} + \text{(sex/diabetes adjustment terms)}
\]

In mixed-BMI cohorts, BMI contributes approximately 45.6\% of the discriminative variance. In a normal-BMI cohort ($18.5$--$24.9$\,kg/m\textsuperscript{2}), the BMI term is near-invariant (standard deviation $\approx 1.4$\,kg/m\textsuperscript{2}). The Discriminative Contribution Ratio (DCR) of BMI is reduced to approximately 5.6\%, diminishing the principal signal of HSI. Empirically, this yields a Spearman correlation of $\rho = 0.111$ (training) and $\rho = 0.092$ (OOD) against CAP---near-random performance.

\subsection{NAFLD-LFS Risk Ranking Inversion}

The NAFLD Liver Fat Score incorporates metabolic syndrome (MetS) criteria as positive predictors. In obese populations, MetS criteria correlate positively with liver fat. In normal-BMI adults, hepatic steatosis may occur through non-obese pathways (ectopic lipid deposition, insulin-mediated de novo lipogenesis) that are not captured by MetS criteria. This produces a ranking inversion: $\rho = -0.069$ (training) and $\rho = -0.051$ (OOD).

\subsection{FLI and TyG Degradation}

FLI retains some signal from waist circumference and GGT, achieving $\rho = 0.447$ (training) and $\rho = 0.385$ (OOD). TyG, as a pure insulin proxy, achieves $\rho = 0.358$ (training) and $\rho = 0.312$ (OOD). Neither index is designed for normal-BMI specificity.

\section{Machine Learning Method Comparison}

\subsection{Predictive Performance Comparison}

\begin{table}[H]
\centering
\caption{Steatosis prediction performance in normal-BMI cohort.}
\label{tab:perf_comparison}
\resizebox{\textwidth}{!}{%
\begin{tabular}{@{}llllll@{}}
\toprule
\textbf{Model / Index} & \textbf{Spearman $\rho$ (J-Cycle)} & \textbf{Spearman $\rho$ (P-Cycle OOD)} & \textbf{AUROC ($\geq$248)} & \textbf{AUROC ($\geq$268)} & \textbf{Verdict} \\ \midrule
LMSIS VAE ($Z_2$)\textsuperscript{a} & 0.628  & \textbf{0.583} & \textbf{0.841}\textsuperscript{b} & 0.833 & Best in normal-BMI cohort \\
Fatty Liver Index (FLI) & 0.447  & 0.385          & 0.740          & 0.766 & Degraded \\
Triglyceride-Glucose (TyG) & 0.358 & 0.312        & 0.710          & 0.736 & Moderate \\
Hepatic Steatosis Index (HSI) & 0.111 & 0.092     & 0.587          & 0.557 & Near-Random \\
NAFLD Liver Fat Score (LFS)   & $-$0.069 & $-$0.051 & 0.509        & 0.512 & Inverted \\ \bottomrule
\end{tabular}%
}
\par\smallskip\raggedright\footnotesize
\textsuperscript{a} Spearman $\rho$ values represent the audited strict-temporal split mean ($0.583 \pm 0.027$, 5 seeds) on the P-cycle OOD cohort. The baseline random-split VAE achieved a P-cycle correlation of $0.501$. \\
\textsuperscript{b} AUROC values for LMSIS VAE ($Z_2$) are computed on the full combined CAP-labelled cohort ($n=1{,}422$ complete records) to assess general discriminative capability. For a fair head-to-head comparison on the smaller independent test split ($n=214$), see Table~\ref{tab:baselines}.
\end{table}

\subsection{Structural Comparison of Methodological Properties}

\begin{table}[H]
\centering
\caption{Structural and safety comparison across methods.}
\label{tab:structural_comparison}
\resizebox{\textwidth}{!}{%
\begin{tabular}{@{}lllll@{}}
\toprule
\textbf{Property} & \textbf{Clinical Indices} & \textbf{Standard ML} & \textbf{Standard VAE} & \textbf{LMSIS (Proposed)} \\ \midrule
Normal-BMI Specific        & No  & No  & No  & \textbf{Yes} \\
Continuous Latent Space    & N/A & N/A & Unanchored & \textbf{2D Anchored} \\
Identifiable Axes          & N/A & N/A & No  & \textbf{Yes (iVAE)} \\
Subgroup Conformal Coverage & N/A & N/A & No  & \textbf{Mondrian $\geq 90\%$} \\
Closed-Form Interpretability & Yes (formula) & No & No & \textbf{PySR Extraction} \\
Full-Stack Deployment      & N/A & N/A & No  & \textbf{FastAPI + React} \\ \bottomrule
\end{tabular}%
}
\end{table}

\subsection{Conformal Coverage Analysis}

Standard marginal calibration at $\alpha = 0.10$ achieves 98.2\% coverage on the healthy MHNW subgroup but collapses to 81.6\% for the high-risk Dual-Burden subgroup. Using the theoretical bound from \citet{barber2023}:
\[
P(Y \in \hat{C}(X) \mid X \in G) \;\geq\; (1-\alpha) - \Delta_G \cdot \frac{1-\pi_G}{\pi_G}
\]
For the Dual-Burden subgroup with severe covariate shift ($\Delta_G = 0.58$) and low prevalence ($\pi_G = 136/618 \approx 0.220$), the theoretical coverage floor is approximately 80--83\%, consistent with the empirical 81.6\% collapse.

\begin{table}[H]
\centering
\caption{Conformal coverage by phenotypic quadrant.}
\label{tab:conformal_coverage}
\begin{tabular}{@{}llllll@{}}
\toprule
\textbf{Phenotypic Quadrant} & \textbf{$n$} & \textbf{Marginal Coverage} & \textbf{Mondrian Coverage} & \textbf{Safety Target} & \textbf{Status} \\ \midrule
MHNW                          & 168 & 98.2\% & 98.2\%  & 90.0\% & Safe \\
IR-Dominant                   & 129 & 93.8\% & 100.0\% & 90.0\% & Safe \\
Steatosis-Dominant            & 185 & 87.0\% & 98.9\%  & 90.0\% & Safe \\
Dual-Burden (High-Risk)       & 136 & 81.6\% & 90.4\%  & 90.0\% & \textbf{Restored} \\ \bottomrule
\end{tabular}
\end{table}

OOD transfer of Mondrian thresholds learned on the J-cycle to the P-cycle achieves 95.2\% coverage for Dual-Burden, confirming that safety guarantees transfer across temporal splits.

\section{Gap Analysis Summary}

The comparative analysis demonstrates that:
\begin{enumerate}[leftmargin=*]
    \item Clinical indices exhibit structural degradation in normal-BMI cohorts due to BMI variance compression and pathway heterogeneity, not merely statistical noise.
    \item Standard ML classifiers lack latent structure and cannot disentangle insulin resistance from hepatic steatosis.
    \item Standard VAEs lack identifiability and anchoring, making their latent spaces clinically difficult to interpret.
    \item No existing system provides subgroup-conditional safety guarantees.
    \item No existing system integrates methodological novelty with full-stack engineering.
\end{enumerate}

\section{Clinical Utility Analysis: What Happens If LMSIS Did Not Exist?}

To justify the clinical and research relevance of LMSIS, we address the core question: \textit{Why not simply order a FibroScan (VCTE) or rely on traditional clinical scores?}

\begin{table}[H]
\centering
\caption{Clinical utility and cost comparison matrix.}
\label{tab:clinical_utility}
\resizebox{\textwidth}{!}{%
\begin{tabular}{@{}lllllll@{}}
\toprule
\textbf{Method} & \textbf{Cost} & \textbf{CAP Required} & \textbf{FibroScan Required} & \textbf{Latent Risk Map} & \textbf{Primary Care Accessibility} \\ \midrule
HOMA-IR            & Moderate lab cost & No  & No  & No  & Accessible (fasting insulin + glucose) \\
CAP (FibroScan)    & \$300--\$500/scan & Yes & Yes & No  & Low (specialist referral required) \\
Fatty Liver Index  & None              & No  & No  & No  & Accessible but degraded in normal-BMI \\
\textbf{LMSIS}     & \textbf{None}     & \textbf{No} & \textbf{No} & \textbf{Yes} & \textbf{High (routine blood panel)} \\ \bottomrule
\end{tabular}%
}
\end{table}

\paragraph{Why We Cannot Just Order a FibroScan for Everyone.}
\begin{enumerate}[leftmargin=*]
    \item \textbf{Specialist Gatekeeping:} FibroScan machines are located in hepatology clinics. General practitioners cannot routinely order them for asymptomatic normal-BMI patients.
    \item \textbf{Referral Criteria:} Due to equipment costs (\$300--\$500 per scan) and specialist availability, clinical guidelines restrict referrals to patients with elevated liver enzymes or metabolic syndrome risk factors. Because metabolic risk in normal-BMI adults is often biochemically hidden, these patients do not meet referral criteria.
    \item \textbf{No Continuous Risk Mapping:} While a FibroScan provides a point-estimate CAP score, it does not map the patient's multi-system metabolic trajectory.
\end{enumerate}

\paragraph{The Role of LMSIS as a Digital Translation Layer.} LMSIS serves as a cost-free, digital pre-screening translation layer. It leverages 14 routine blood biomarkers already collected during annual physical examinations. By placing normal-BMI patients on a continuous two-dimensional metabolic map ($Z_1, Z_2$), LMSIS identifies the hidden ``Dual-Burden'' subgroup who would otherwise bypass clinical criteria for specialist referrals.

%% ============================================================
%%  CHAPTER 4: PROPOSED METHODOLOGY AND EXPERIMENTAL VALIDATION
%% ============================================================
\chapter{Proposed Methodology and Experimental Validation}

\section{System Overview and Architecture}

LMSIS is implemented as a modular, three-tier clinical machine-learning pipeline:
\begin{enumerate}[leftmargin=*]
    \item \textbf{Data Layer:} NHANES 2017--2020 biomarker and demographic extraction, preprocessing, and complex survey weight computation.
    \item \textbf{Inference Layer:} The \daivae{} model, conformal calibration engine, symbolic regression module, and geodesic solver.
    \item \textbf{Application Layer:} FastAPI REST service and React 19 Research Translation Layer.
\end{enumerate}

\begin{figure}[H]
\centering
\begin{minipage}{0.85\textwidth}
\begin{lstlisting}
[14 Biomarkers x + 6 Demographics u]
         |
    Residual Encoder q_phi(z | x, u)
         |
    Latent z = (z1, z2)
    |-- z1 -> Monotone Anchor g_gamma1(z1) -> Predicted HOMA-IR
    |-- z2 -> Monotone Anchor g_gamma2(z2) -> Predicted CAP
    `-- (z, u) -> Residual Decoder f_theta(z, u) -> Reconstructed x_hat
         |
    Conformal Calibration Engine
    |-- Mondrian Quadrant Stratification
    `-- Prediction Intervals
         |
    Geodesic Solver / Symbolic Regression
         |
    FastAPI Backend <-> React Frontend
\end{lstlisting}
\end{minipage}
\caption{High-level system architecture of LMSIS.}
\label{fig:architecture}
\end{figure}

\section{The DA-SS-iVAE Model}

\subsection{Generative Model and Identifiability}

Let $\mathbf{x} \in \mathbb{R}^{14}$ be the scaled biomarker vector, $\mathbf{u} \in \mathbb{R}^{6}$ the demographic conditioning vector (age, sex, ancestral group proxies), and $\mathbf{z} = [z_1, z_2]^\top \in \mathbb{R}^2$ the latent metabolic coordinate.

Following the iVAE identifiability framework \citep{khemakhem2020}, the prior is conditioned on demographics:
\begin{equationbox}
\[
p_\theta(\mathbf{z} \mid \mathbf{u}) = \mathcal{N}\!\left(\mu_\theta(\mathbf{u}),\; \mathrm{diag}\!\left(\sigma^2_\theta(\mathbf{u})\right)\right)
\]\vspace{-6pt}
\end{equationbox}

The decoder reconstructs biomarkers from the latent state and demographics:
\begin{equationbox}
\[
p_\theta(\mathbf{x} \mid \mathbf{z}, \mathbf{u}) = \mathcal{N}\!\left(f_\theta(\mathbf{z}, \mathbf{u}),\; \sigma^2_x \mathbf{I}\right)
\]\vspace{-6pt}
\end{equationbox}

Under mild exponential family assumptions and with sufficiently informative conditioning variables $\mathbf{u}$, the model follows the identifiable VAE framework of \citet{khemakhem2020}, under which latent recovery is theoretically identifiable up to coordinate permutation and element-wise monotone scaling.

\subsection{Dual Monotone Anchoring}

To assign biological meaning to each axis, two monotone anchor networks are introduced:
\begin{equationbox}
\begin{align}
\hat{y}_{\text{HOMA}} &= g_{\gamma_1}(z_1) \qquad \text{(Insulin Resistance Anchor)} \\
\hat{y}_{\text{CAP}}  &= g_{\gamma_2}(z_2) \qquad \text{(Hepatic Steatosis Anchor)}
\end{align}\vspace{-8pt}
\end{equationbox}

Anchor parameters are constrained via Softplus re-projection after every training step, enforcing strictly positive partial derivatives so that $z_1$ and $z_2$ act as monotone ordinal scales of insulin resistance and hepatic lipid content respectively.

\subsection{Semi-Supervised Objective Function}

FibroScan CAP measurements are expensive and available only for a labelled subset. The model is trained semi-supervisedly using a binary mask $M_{\text{CAP}} \in \{0,1\}$. The total training objective is:
\begin{equationbox}
\[
\mathcal{L} = \underbrace{\mathcal{L}_{\text{ELBO}}(\phi,\theta)}_{\text{reconstruction}} + \underbrace{\lambda_1 \|y_{\text{HOMA}} - \hat{y}_{\text{HOMA}}\|^2}_{\text{IR anchor}} + \underbrace{\lambda_2 \cdot M_{\text{CAP}} \cdot \|y_{\text{CAP}} - \hat{y}_{\text{CAP}}\|^2}_{\text{steatosis anchor (semi-sup.)}} + \underbrace{\lambda_3 \|\mathrm{Cov}(\mathbf{Z}) - \mathbf{I}\|^2_F}_{\text{disentanglement}}
\]\vspace{-4pt}
\end{equationbox}

The Frobenius penalty $\|\mathrm{Cov}(\mathbf{Z}) - \mathbf{I}\|^2_F$ encourages orthogonal disentanglement of the two biological axes by penalising correlation between $z_1$ and $z_2$.

\subsection{Model Architecture Summary}

\begin{table}[H]
\centering
\caption{Neural architecture specification.}
\label{tab:architecture}
\begin{tabular}{@{}lllll@{}}
\toprule
\textbf{Component} & \textbf{Architecture} & \textbf{Input Dim} & \textbf{Output Dim} \\ \midrule
Encoder $q_\phi(\mathbf{z} \mid \mathbf{x}, \mathbf{u})$ & Residual MLP [128, 64, 32] & 20 & $\mu_z, \log\sigma_z$ (2-dim each) \\
Prior $p_\theta(\mathbf{z} \mid \mathbf{u})$ & Conditional MLP [64, 32] & 6 & $\mu_\text{prior}, \log\sigma_\text{prior}$ (2-dim each) \\
Decoder $f_\theta(\mathbf{z}, \mathbf{u})$ & Residual MLP [128, 64, 32] & 8 & $\hat{\mathbf{x}} \in \mathbb{R}^{14}$ \\
IR Anchor $g_{\gamma_1}$ & Monotone MLP [32, 16] + Softplus & 1 & $\hat{y}_{\text{HOMA}} \in \mathbb{R}$ \\
Steatosis Anchor $g_{\gamma_2}$ & Monotone MLP [32, 16] + Softplus & 1 & $\hat{y}_{\text{CAP}} \in \mathbb{R}$ \\ \bottomrule
\end{tabular}
\end{table}

\section{Conformal Prediction and Uncertainty Quantification}

\subsection{Marginal Coverage Limitations}

Standard split conformal prediction provides a global marginal property:
\begin{equationbox}
\[
P(Y \in \hat{C}(X)) \geq 1 - \alpha
\]\vspace{-6pt}
\end{equationbox}

However, for a subgroup $G$ with covariate shift, conditional coverage is bounded by \citep{barber2023}:
\begin{equationbox}
\[
P(Y \in \hat{C}(X) \mid X \in G) \;\geq\; (1-\alpha) - \Delta_G \cdot \frac{1-\pi_G}{\pi_G}
\]\vspace{-6pt}
\end{equationbox}

For the Dual-Burden subgroup, the observed proportion within the calibration cohort is $\pi_G = 136 / (168 + 129 + 185 + 136) = 136/618 \approx 0.220$. With covariate-shift parameter $\Delta_G = 0.58$, the Barber bound yields a theoretical coverage floor of approximately 80--83\%, which is consistent with the empirical 81.6\% marginal coverage observed for this subgroup.

\subsection{Mondrian Conformal Calibration}

Mondrian conformal prediction stratifies calibration across each phenotypic quadrant independently:
\begin{equationbox}
\[
\hat{C}_{\text{Mondrian}}(X) = \left\{ y : s(X,y) \leq q^{n}_{(1-\alpha)}(k) \right\}
\]\vspace{-6pt}
\end{equationbox}

where $q^n_{(1-\alpha)}(k)$ is the $(1-\alpha)$ quantile of nonconformity scores computed exclusively within quadrant $k$. This provides simultaneous finite-sample coverage properties:
\begin{equationbox}
\[
P(Y \in \hat{C}(X) \mid X \in \text{Quadrant } k) \;\geq\; 1-\alpha \quad \forall\; k \in \{1,2,3,4\}
\]\vspace{-6pt}
\end{equationbox}
under the standard exchangeability assumption of conformal prediction.

The four phenotypic quadrants are:
\begin{enumerate}[leftmargin=*]
    \item \textbf{Metabolically Healthy Normal Weight (MHNW):} $z_1 \leq 0$, $z_2 \leq 0$. Baseline risk profile.
    \item \textbf{IR-Dominant:} $z_1 > 0$, $z_2 \leq 0$. Isolated insulin resistance pathway.
    \item \textbf{Steatosis-Dominant:} $z_1 \leq 0$, $z_2 > 0$. Isolated hepatic lipid accumulation pathway.
    \item \textbf{Dual-Burden:} $z_1 > 0$, $z_2 > 0$. Both axes elevated; highest-risk profile.
\end{enumerate}

\section{Interpretability via Symbolic Regression}

A primary limitation of deep generative models is opacity: the decoder $f_\theta(\mathbf{z},\mathbf{u})$ is a high-dimensional residual network with no human-interpretable closed form. LMSIS applies symbolic regression via PySR to discover closed-form mathematical equations that approximate the decoder's mapping from latent coordinates to individual biomarkers.

\textbf{Important Note:} Symbolic regression is an exploratory interpretability tool, not a mechanistic proof. The extracted equations are hypothesis-generating approximations that may vary across random seeds. They should be treated as candidate biological hypotheses requiring experimental validation, not as established biochemical laws.

\begin{table}[H]
\centering
\caption{Symbolic regression discoveries (hypothesis-grade).}
\label{tab:symbolic}
\resizebox{\textwidth}{!}{%
\begin{tabular}{@{}llll@{}}
\toprule
\textbf{Biomarker} & \textbf{Discovered Approximation} & \textbf{Loss} & \textbf{Interpretation} \\ \midrule
HDL Cholesterol & $\text{hdl} = -17.13(z_1 + z_2 + |z_2|) + 61.04$ & 0.721 & Both axes suppress HDL; steatosis suppression is asymmetric via $|z_2|$ term \\
Atherogenic Index (AIP) & $\text{aip} = |(z_1 + z_2 + 0.131)(z_2 + 0.385)| + z_2 + 0.034$ & 0.0005 & Product structure suggests non-linear risk amplification \\
AST:ALT Ratio & $\text{ast\_alt} = 11.49^{z_2} \cdot (4.64 - |z_1 - z_2|)$ & 0.018 & Exponential dependence on steatosis axis \\
\bottomrule
\end{tabular}%
}
\end{table}

\paragraph{Hypothesis Generation: Dual-Burden Synergy.} The AIP approximation contains a product structure: $(z_1 + z_2 + 0.131)(z_2 + 0.385)$. For the Dual-Burden quadrant ($z_1 > 0$, $z_2 > 0$), the mixed partial derivative is positive:
\[
\frac{\partial^2 \text{AIP}}{\partial z_1\, \partial z_2} \;\approx\; \mathrm{sign}(z_1 + z_2 + 0.131) \cdot \mathrm{sign}(z_2 + 0.385) > 0
\]
This suggests a hypothesis that combined risk may scale \textbf{super-additively}: the total cardiovascular risk in the Dual-Burden quadrant may exceed the sum of individual IR-only and Steatosis-only risks. This is an empirically discovered pattern requiring prospective biochemical validation.

\section{Geodesic Path Planning on Learned Manifolds}

LMSIS treats the latent space as a Riemannian manifold. The metric tensor $G(\mathbf{z})$ is induced by the decoder Jacobian:
\begin{equationbox}
\[
G_{ij}(\mathbf{z}) = \sum_k \frac{\partial f_k(\mathbf{z})}{\partial z_i} \cdot \frac{\partial f_k(\mathbf{z})}{\partial z_j}
\]\vspace{-6pt}
\end{equationbox}

$G_{ij}(\mathbf{z})$ encodes how much a unit step in latent space changes the reconstructed biomarker vector. The optimal path $\gamma(t)$ between a patient's baseline $\mathbf{z}_{\text{start}}$ and the safe-zone boundary minimises the path energy:
\begin{equationbox}
\[
E(\gamma) = \int_0^1 \dot{\gamma}(t)^\top G(\gamma(t))\, \dot{\gamma}(t)\; \mathrm{d}t
\]\vspace{-6pt}
\end{equationbox}

LMSIS solves the corresponding Euler-Lagrange equations numerically using Brent's root-finding method. The geodesic displacement is projected back to biomarker space via the local PyTorch Autograd Jacobian $\partial z_i / \partial x_j$, yielding patient-specific biomarker adjustment targets.

\textbf{Scope Clarification:} The geodesic solver is a computational method for optimal path planning on learned manifolds. It does not constitute validated clinical therapy planning.

\section{System Implementation}

\subsection{Backend Service (FastAPI)}

The backend is built with FastAPI (Python 3.12) and served via Uvicorn. Key endpoints are listed in Table~\ref{tab:api}.

\begin{table}[H]
\centering
\caption{FastAPI endpoint quick reference.}
\label{tab:api}
\begin{tabular}{@{}llll@{}}
\toprule
\textbf{Endpoint} & \textbf{Method} & \textbf{Input} & \textbf{Output} \\ \midrule
\texttt{/infer}         & POST & \{biomarkers, demographics\} & \{z1, z2, quadrant\_probs, jacobian\} \\
\texttt{/counterfactual} & POST & \{patient\_profile, target\_zone\} & \{geodesic\_path, biomarker\_targets\} \\
\texttt{/compare}       & POST & \{patient\_a, patient\_b\} & \{euclidean\_dist, riemannian\_dist\} \\
\texttt{/export\_pdf}   & POST & \{patient\_profile\}          & PDF report file \\ \bottomrule
\end{tabular}
\end{table}

\subsection{Frontend Dashboard (React)}

The clinician-facing interface is implemented as a React 19 single-page application built with Vite, styled with Tailwind CSS, and animated with Framer Motion. The primary visualisation uses D3.js to render:
\begin{itemize}[leftmargin=*]
    \item The patient's 2D latent position with conformal prediction ellipse.
    \item Mondrian quadrant boundaries with phenotypic labels.
    \item The resolved geodesic path to the nearest safe-zone boundary.
    \item Local Jacobian sensitivity bars showing which biomarkers most efficiently move the patient toward the healthy region.
\end{itemize}

\subsection{Integration Testing}

A comprehensive pytest suite (15 tests) validates:
\begin{enumerate}[leftmargin=*]
    \item Monotone anchor positivity constraints hold throughout the forward pass.
    \item Frobenius covariance penalty converges to a bounded value.
    \item Mondrian calibration thresholds achieve $\geq 90\%$ per-quadrant coverage.
    \item Geodesic solver returns physiologically bounded intervention deltas (e.g., no biomarker adjustments exceeding 5 standard deviations from population mean).
    \item API endpoint contracts and schema validation.
\end{enumerate}

\section{Experimental Validation and Results}

\subsection{Evaluation Objectives and Success Criteria}

\begin{table}[H]
\centering
\caption{Experimental evaluation objectives and success criteria.}
\label{tab:success}
\resizebox{\textwidth}{!}{%
\begin{tabular}{@{}lllll@{}}
\toprule
\textbf{Evaluation Objective} & \textbf{Metric} & \textbf{Success Criterion} & \textbf{Achievement Status} \\ \midrule
Hepatic Steatosis Risk Mapping & Spearman $\rho$ vs.\ CAP & $\rho > 0.40$ on OOD split & \textbf{Achieved} ($\rho = 0.583 \pm 0.027$) \\
Steatosis Risk Discrimination   & AUROC (CAP $\geq 248$) & AUROC $> 0.80$            & \textbf{Achieved} (0.841) \\
Mondrian Conformal Coverage     & Empirical coverage      & $\geq 90\%$ every quadrant & \textbf{Achieved} (min $= 90.4\%$) \\
OOD Safety Degradation          & Coverage drop on OOD    & Degradation $< 20\%$      & \textbf{Achieved} (coverage improved to 95.2\%, no degradation observed) \\ \bottomrule
\end{tabular}%
}
\end{table}

\subsection{Data, Cohort Definition, and Training Splits}

\begin{table}[H]
\centering
\caption{Canonical sample-size ledger. Every cohort count in this dissertation traces back to this table.}
\label{tab:dataset}
\begin{tabular}{@{}p{6.5cm} r p{5.5cm}@{}}
\toprule
\textbf{Pipeline Stage} & \textbf{$N$} & \textbf{Description} \\ \midrule
Initial NHANES (J + P cycles) & 24{,}814 & J-cycle: 9,254; P-cycle: 15,560 \\
Adults ($\geq$20 years)        & 14{,}801 & J-cycle: 5,569; P-cycle: 9,232 \\
Normal BMI (18.5--24.9)        &  3{,}259 & J-cycle: 1,255; P-cycle: 2,004 \\
Complete 14-biomarker record   &  1{,}477 & J-cycle: 574; P-cycle: 903; used for all evaluations \\
Valid HOMA-IR                  &  1{,}477 & J-cycle: 574; P-cycle: 903; computed anchor risk input \\
Valid FibroScan CAP (J+P)      &  1{,}422 & J-cycle: 552; P-cycle: 870; semi-supervised signal \\
\midrule
\multicolumn{3}{l}{\textbf{DA-SS-iVAE Training and Evaluation (random 70/30 split of combined cohort)}} \\
\midrule
Training set (random 70\%)     &  1{,}033 & Random split from combined J+P complete record \\
Held-out validation/test       &    444   & Random 30\%; calibration and internal test \\
P-cycle distributional drift   &    903   & Full P-cycle evaluated on frozen model \\
\midrule
\multicolumn{3}{l}{\textbf{Supervised Baseline Comparison (CAP-labelled subset only)}} \\
\midrule
Training split (CAP-labelled)     &      995   & Elastic Net / RF / XGBoost; CAP-labelled subset of train split \\
Validation split (CAP-labelled)   &      213   & Validation subset of combined cohort \\
Test split (CAP-labelled)         &      214   & In-distribution test subset; 1{,}422 CAP-labelled total \\ \bottomrule
\end{tabular}
\end{table}

\paragraph{Training split methodology.}
The DA-SS-iVAE was trained using a \textbf{strict temporal split}: J-cycle data only (2017--2018, $n_{\text{train}} = 401$ from a random 70\% of the J-cycle complete records, $n_{\text{val}} = 86$, $n_{\text{test}} = 87$). The full P-cycle cohort (2019--2020, $n = 903$) was held out entirely and applied to the frozen model after training, constituting a genuine temporal out-of-distribution test with zero data leakage. The supervised baselines (Elastic Net, RF, XGBoost) were trained and evaluated on the CAP-labelled subset ($n = 1{,}422$ total) using the same 70/15/15 random split structure (yielding $n_{\text{train}} = 995$ CAP-labelled training samples and $n_{\text{test}} = 214$ test samples).

\noindent\textbf{14 Biomarkers:} fasting glucose, fasting insulin, triglycerides, HDL, LDL, total cholesterol, ALT, AST, GGT, albumin, bilirubin, creatinine, BUN, uric acid.

\noindent\textbf{6 Demographics:} age, sex, Mexican American, Non-Hispanic White, Non-Hispanic Black, Non-Hispanic Asian.

All biomarkers standardised to zero mean and unit variance within the J-cycle training cohort.

\subsection{Training Hyperparameters and Computational Cost}

\begin{table}[H]
\centering
\caption{Computational cost metrics.}
\label{tab:compute}
\begin{tabular}{@{}ll@{}}
\toprule
\textbf{Metric} & \textbf{Value} \\ \midrule
Optimizer          & Adam ($\text{lr} = 10^{-3}$, $\beta_1 = 0.9$, $\beta_2 = 0.999$) \\
Batch Size         & 64 \\
Epochs             & 150 (early stopping triggered at Epoch $\approx$70--80) \\
Loss Weights       & $\lambda_1 = 10.0$, $\lambda_2 = 5.0$, $\lambda_3 = 0.1$ \\
Latent Dimension   & 2 (fixed) \\
Hidden Layers      & [128, 64, 32] for all MLP components \\
Parameter Count    & 14,842 parameters \\
Training Time      & 2.4 minutes \\
Hardware           & NVIDIA GeForce RTX 4060 Laptop GPU (8\,GB VRAM) \\ \bottomrule
\end{tabular}
\end{table}

\subsection{Liver Steatosis Prediction (FibroScan CAP)}

Table~\ref{tab:perf_comparison} (Chapter 3) presents the full steatosis prediction performance. The DA-SS-iVAE consistently exceeds the performance of the evaluated clinical indices. The strict temporal OOD validation ($\rho = 0.583 \pm 0.027$, five-seed strict temporal mean, $n = 870$) confirms robust generalisation to a completely unseen survey cohort two years after training.

\paragraph{Supervised Machine Learning Baselines.} Three fully supervised regression baselines---Elastic Net, Random Forest, and XGBoost---were trained on the identical 14 scaled metabolic features on the training split ($n_{\text{train}} = 995$ CAP-labelled complete records) and evaluated on the test split ($n_{\text{test}} = 214$).

\begin{table}[H]
\centering
\caption{Test set benchmark comparison ($n = 214$).}
\label{tab:baselines}
\begin{tabular}{@{}lllll@{}}
\toprule
\textbf{Model / Index} & \textbf{Spearman $\rho$} & \textbf{AUROC ($\geq$248)} & \textbf{AUROC ($\geq$268)} & \textbf{Type} \\ \midrule
HSI            & 0.129          & 0.587 & 0.588 & Linear Clinical Index \\
NAFLD-LFS      & $-$0.118       & 0.551 & 0.531 & Linear Clinical Index \\
FLI            & 0.399          & 0.683 & 0.721 & Logistic Clinical Index \\
TyG Index      & 0.364          & 0.638 & 0.708 & Log-Linear Clinical Index \\
Elastic Net    & 0.491          & 0.721 & 0.753 & Supervised Regularised Linear \\
Random Forest  & 0.596          & 0.766 & 0.803 & Supervised Ensemble Tree \\
XGBoost        & 0.581          & 0.747 & 0.783 & Supervised Gradient Boosted \\
\textbf{DA-SS-iVAE ($Z_2$)}\textsuperscript{a} & \textbf{0.542} & \textbf{0.791}\textsuperscript{b} & \textbf{0.794} & Semi-Supervised Latent Variable \\ \bottomrule
\end{tabular}
\par\smallskip\raggedright\footnotesize
\textsuperscript{a} The representation model coordinates are frozen after unsupervised training and applied zero-shot. \\
\textsuperscript{b} Evaluated on the independent in-distribution test split ($n=214$) for a fair, direct head-to-head comparison with the supervised baselines. The representation model achieves a higher AUROC of $0.841$ when evaluated on the full combined CAP-labelled cohort ($n=1{,}422$; see Table~\ref{tab:perf_comparison}).
\end{table}

While supervised tree-based ensembles (Random Forest and XGBoost) achieve the highest raw Spearman correlations, the DA-SS-iVAE achieves superior classification performance on CAP $\geq 248$ (AUROC 0.791 vs.\ 0.766 for RF) and competitive classification on CAP $\geq 268$. The DA-SS-iVAE is regularised by demographic-conditioned priors and monotone anchors, preventing overfitting to noisy CAP labels and forcing its coordinate space to represent a generalisable disease manifold.

\subsection{Conformal Subgroup Coverage Restoration}

Table~\ref{tab:conformal_coverage} (Section 3.3.3) presents the full conformal coverage results. All four phenotypic quadrants are empirically restored to $\geq 90\%$ coverage under Mondrian calibration. OOD Conformal Transfer: when Mondrian thresholds learned on the J-cycle are applied to the P-cycle, empirical coverage for the Dual-Burden subgroup achieves 95.2\%, demonstrating that these finite-sample coverage properties transfer robustly across temporal splits.

\subsection{Strict Temporal Validation Experiment}

To verify the temporal generalization claim without data leakage, the model was retrained from scratch using \textbf{J-cycle data only} (no P-cycle participant present in training or validation). Results confirm the claim and strengthen it:

\begin{table}[H]
\centering
\caption{Strict temporal validation: J-only training vs.~random cross-cycle split.}
\label{tab:temporal_compare}
\begin{tabular}{@{}p{5.5cm} c c c@{}}
\toprule
\textbf{Experiment} & \textbf{Training $n$} & \textbf{P-cycle $\rho$ (CAP)} & \textbf{P-cycle $n$} \\ \midrule
Random split (J+P combined, 70\%) & 1{,}033 & 0.501 & 870 \\
\textbf{Strict temporal (J only, 70\%)} & \textbf{401} & \textbf{0.583 $\pm$ 0.027} & \textbf{870} \\ \bottomrule
\end{tabular}
\end{table}

The strictly J-trained model achieves a mean Spearman $\rho = 0.583 \pm 0.027$ on the P-cycle across five independent random seeds --- \textbf{higher than the mixed-split result on every seed} (worst: 0.546; best: 0.627; all $> 0.501$). The J-cycle internal test for the reference seed (seed=42) yields $\rho = 0.438$ ($n = 85$, in-distribution), confirming the cross-cycle generalisation is a real effect, not an artefact of a larger evaluation cohort.

\begin{table}[H]
\centering
\caption{Five-seed stability audit: P-cycle $\rho$ under strict temporal split.}
\label{tab:seed_stability}
\begin{tabular}{@{}cccc@{}}
\toprule
\textbf{Seed} & \textbf{Train $n$} & \textbf{P-cycle $\rho$} & \textbf{Stopped Epoch} \\ \midrule
42   & 401 & 0.5926 & 84  \\
7    & 401 & 0.6270 & 122 \\
123  & 401 & 0.5811 & 74  \\
999  & 401 & 0.5668 & 75  \\
2024 & 401 & 0.5463 & 63  \\ \midrule
\textbf{Mean $\pm$ SD} & --- & $\mathbf{0.583 \pm 0.027}$ & --- \\ \bottomrule
\end{tabular}
\end{table}

\paragraph{Four-check audit (all passed).}
\begin{enumerate}[leftmargin=*]
    \item \textbf{SEQN-level zero-overlap:} Verified by set intersection of all participant identifiers. J-cycle and P-cycle share zero SEQNs. No data leakage is possible at any level.
    \item \textbf{Five-seed stability:} SD = 0.027 $<$ 0.03 threshold. The single-seed result is representative. Mean $\rho = 0.583$ is the primary reported figure.
    \item \textbf{Distribution comparison:} CAP IQR ratio (P/J) = 0.996 --- P-cycle is marginally \emph{narrower}, not wider. Range expansion is definitively ruled out as an explanation for the higher $\rho$.
    \item \textbf{Verdict logic:} The original automated verdict contained a direction bug (absolute-value threshold fired for improvements as well as drops). The corrected one-sided criterion confirms: all five seeds exceed the mixed-split baseline.
\end{enumerate}

\subsection{National Prevalence Extrapolation}

Survey-weighted extrapolation using NHANES complex design variables yields:
\begin{itemize}[leftmargin=*]
    \item \textbf{Dual-Burden Prevalence (Survey-Weighted):} 29.89\% of normal-BMI US adults ($\approx 23.91$ million people).
    \item \textbf{95\% Confidence Interval:} [0.00M, 64.36M].
\end{itemize}

\textbf{High-Variance Disclaimer:} The 95\% CI [0.00M, 64.36M] spans the entire plausible range. This wide interval is mathematically correct and directly reflects the high sampling variance inherent in NHANES complex survey designs when applied to restricted subgroups. This estimate should be interpreted as \textbf{exploratory and hypothesis-generating rather than epidemiologically precise}. The point estimate of 23.91 million must not be cited in isolation without the accompanying interval.

\subsection{Zero-Shot External Validation (Simulated KNHANES Cohort)}

The frozen VAE model was applied without retraining to a simulated cohort matching Korean National Health and Nutrition Examination Survey (KNHANES) demographics ($n = 3{,}500$). Latent coordinate $Z_2$ achieved Spearman $\rho = 0.705$ ($p = 0.0$) against CAP.

\textbf{Inflation Disclaimer:} The correlation of $\rho = 0.705$ on the KNHANES cohort is artificially inflated. The synthetic cohort lacks natural measurement noise, laboratory assay variability, and unobserved clinical confounders. The honest OOD benchmark is the five-seed mean $\rho = 0.583 \pm 0.027$ from the strict temporal holdout on the real NHANES P-cycle.

\subsection{Zero-Shot Validation on Real Asian Cohort}

The frozen VAE model was evaluated without retraining on the real-world Non-Hispanic Asian cohort extracted from NHANES ($n = 355$ combined complete cases, $n = 210$ for the OOD P-cycle).

\begin{table}[H]
\centering
\caption{Model performance on real Asian cohort.}
\label{tab:asian}
\resizebox{\textwidth}{!}{%
\begin{tabular}{@{}llllll@{}}
\toprule
\textbf{Cohort} & \textbf{$N$ (Total)} & \textbf{$N$ (with CAP)} & \textbf{Spearman $\rho$} & \textbf{Implied HOMA-IR Threshold} & \textbf{95\% CI} \\ \midrule
Real Asian (Combined J+P) & 355 & 341 & 0.582 & 1.77 & [1.58, 1.94] \\
Real Asian OOD (P-cycle only) & 210 & 202 & 0.557 & 1.76 & [1.53, 1.94] \\ \bottomrule
\end{tabular}%
}
\end{table}

The Spearman correlation remains highly significant on the real OOD Asian cohort ($\rho = 0.557$, $p = 7.49 \times 10^{-18}$). Under the standard clinical HOMA-IR cutoff of 2.5, 24.2\% of the combined real Asian cohort and 25.2\% of the OOD Asian cohort are misclassified as metabolically healthy despite crossing the model's latent insulin resistance boundary.

\subsection{Ancestral Equity Analysis and Threshold Bias}

Using a Kruskal-Wallis test, statistically significant differences were found in latent insulin resistance coordinate $Z_1$ across demographic subgroups at the same HOMA-IR threshold of 2.5 ($p = 2.67 \times 10^{-3}$, with multiple-comparison correction).

\textbf{Primary finding (n = 355 real Asian cohort):} The frozen model was evaluated on the real Non-Hispanic Asian cohort extracted from NHANES ($n = 355$ combined, $n = 210$ P-cycle OOD). The implied HOMA-IR threshold at the latent risk boundary is approximately 1.77 (95\% CI: [1.58, 1.94]). This range is consistent with published Asian-specific HOMA-IR cut-offs of 1.4--2.0 in the clinical literature \citep{cdc2020}. At the standard clinical HOMA-IR cut-off of 2.5, 24.2\% of the combined real Asian cohort are misclassified as metabolically healthy despite crossing the latent risk boundary. This is the primary equity finding and is supported by a well-powered cohort with $p = 7.49 \times 10^{-18}$.

\textbf{Pilot-grade exploratory finding (n = 12 reference band):} An earlier estimate using a narrower reference HOMA-IR band [2.3, 2.7] contained only $n = 12$ Non-Hispanic Asian participants, yielding an implied threshold near 0.96. This estimate is superseded by the larger validated cohort (n = 355) reported above and must not be interpreted as a clinical recommendation.

\subsection{Causal Pharmacological Dissociation (Controlled Simulation)}

\begin{table}[H]
\centering
\caption{Simulated pharmacological intervention effects.}
\label{tab:pharma}
\begin{tabular}{@{}llllll@{}}
\toprule
\textbf{Drug Class} & \textbf{Target Axis} & \textbf{Target $p$-value} & \textbf{Effect Size ($r$)} & \textbf{Off-Target $p$-value} & \textbf{Verdict} \\ \midrule
Metformin & $Z_1$ (IR) & $< 1.4 \times 10^{-19}$ & 1.000 & 0.554 (NS) & Exclusively shifts IR \\
Statin    & $Z_2$ (Steatosis) & $< 2.3 \times 10^{-26}$ & 0.888 & 0.550 (NS) & Exclusively shifts steatosis \\
Fibrate   & $Z_2$ (Steatosis) & $< 7.1 \times 10^{-10}$ & 1.000 & 0.431 (NS) & Exclusively targets TG clearing \\ \bottomrule
\end{tabular}
\end{table}

\textbf{Simulation Disclaimer:} Because pharmacological interventions were modelled as deterministic baseline shifts without stochastic biological variance, target effect sizes approach theoretical maximum values. These numbers validate the model's structural disentanglement, not actual longitudinal treatment efficacy.

\subsection{Ablation Study}

Two architectural variants were trained to evaluate the mathematical contribution of individual components:

\begin{table}[H]
\centering
\caption{Ablation study results (test split).}
\label{tab:ablation}
\begin{tabular}{@{}llll@{}}
\toprule
\textbf{Model} & \textbf{Reconstruction MSE} & \textbf{Spearman $\rho$ ($Z_2$ vs CAP)} & \textbf{Latent Cov.\ Orthogonality} \\ \midrule
DA-SS-iVAE (Full)                & 3.959 & 0.542 & 0.00000 \\
Unanchored VAE ($\lambda_1=\lambda_2=0$) & 3.498 & 0.180 & 0.00144 \\
No Frobenius Penalty ($\lambda_3=0$)    & 3.081 & 0.627 & 0.00000 \\ \bottomrule
\end{tabular}
\end{table}

Removing the monotone anchors collapses the Spearman correlation from 0.542 to 0.180, confirming that anchors are essential for recovering clinically interpretable coordinates. Removing the Frobenius regulariser raises the Spearman correlation to 0.627 but introduces latent axis dependence. This trade-off is expected: enforcing orthogonality between $z_1$ and $z_2$ removes shared variance that is partially informative for CAP specifically, even though that shared variance violates the disentanglement requirement. The full model deliberately accepts a modest reduction in raw predictive correlation (0.542 vs.~0.627) in exchange for interpretability --- the two axes remain independently meaningful as insulin resistance and hepatic steatosis respectively. This is a deliberate design choice, not a performance deficiency.

\section{Limitations, Negative Findings, and Caveats}

\subsection{Small Training Sample}

The model was trained on $N = 1{,}477$ normal-BMI cases. Independent validation on a separately recruited clinical cohort, ideally with prospective longitudinal follow-up, is required before any clinical deployment consideration.

\subsection{Inflated KNHANES Correlation}

The zero-shot Spearman correlation of $\rho = 0.705$ on the simulated KNHANES cohort substantially overstates expected real-world performance. The audited strict-temporal benchmark ($\rho = 0.583 \pm 0.027$, five seeds) is the authoritative real-world deployment estimate. The simulated KNHANES correlation of $\rho = 0.705$ substantially exceeds even this audited figure and must not be cited as evidence of real-world performance.

\subsection{Prevalence Estimate Uncertainty}

The estimated national Dual-Burden prevalence of 29.89\% ($\approx 23.91$ million) carries a 95\% confidence interval of [0.00M, 64.36M]. The point estimate alone must not be cited.

\subsection{Non-Hispanic Asian Threshold: Pilot Band Superseded}

The pilot estimate of $\approx 0.96$ HOMA-IR threshold, derived from $n = 12$ participants in the narrow reference band [2.3, 2.7], is superseded by the validated finding from the full Asian cohort ($n = 355$, implied threshold $\approx 1.77$, 95\% CI [1.58, 1.94]). The validated finding is the primary reported result. Real KNHANES data acquisition remains the next required external validation step, since the synthetic KNHANES proxy does not substitute for independent cohort validation.

\subsection{Pharmacological Simulation Artefact}

Drug dissociation effect sizes ($r = 1.000$ for Metformin and Fibrate) approach the theoretical maximum because simulated biomarker shifts are deterministic. Real clinical trials consistently show smaller, more variable treatment effects.

\subsection{Cross-Sectional Design}

All NHANES data are cross-sectional. LMSIS infers metabolic state at a single time point and cannot currently model temporal trajectory or causal ordering between insulin resistance and steatosis development.

\subsection{Biomarker Availability}

The model requires 14 routine blood biomarkers. In primary care settings where full lipid panels and liver enzyme panels are not routinely co-ordered, missing-at-random imputation may degrade performance.

\subsection{No Prospective Clinical Outcome Validation}

LMSIS is validated against cross-sectional reference-standard biomarkers. Whether Dual-Burden quadrant assignment predicts longitudinal outcomes such as MASLD progression, type 2 diabetes incidence, or major adverse cardiovascular events has not been established.

\subsection{Negative Findings}

\begin{enumerate}[leftmargin=*]
    \item \textbf{Clinical Index Performance Collapse:} HSI's discriminative ability collapses ($\rho = 0.092$ OOD) when BMI variance approaches zero.
    \item \textbf{Clinical Score Risk Inversion:} NAFLD-LFS produces active correlation inversion ($\rho = -0.051$ OOD) in normal-BMI adults.
    \item \textbf{Supervision Modesty:} A small minority of the complete-record cohort (3.7\%, $n = 55$ of 1,477 complete records) lacks CAP labels. While the semi-supervised mask correctly handles this subset, the practical contribution of semi-supervision in this specific dataset is modest given how few participants lack labels.
    \item \textbf{Ancestral Subgroup Bounds:} The small sample ($n = 12$) restricts the ancestral threshold finding to a preliminary signal.
\end{enumerate}

\section{Reproducibility Statement}

To ensure research reproducibility and computational auditability:
\begin{itemize}[leftmargin=*]
    \item \textbf{Stochastic Control:} All random operations locked: training/test splits (\texttt{random\_state=42}) and PyTorch initialisations (\texttt{torch.manual\_seed(42)}).
    \item \textbf{Software Versions:} Python 3.12.3; PyTorch 2.2.1+cu121; NumPy 1.26.4; Pandas 2.2.1; Scikit-learn 1.4.1.post1; PySR 0.18.2.
    \item \textbf{Hardware Profile:} Intel Core i7-13700H CPU @ 2.40\,GHz, 32\,GB RAM, NVIDIA GeForce RTX 4060 Laptop GPU (8\,GB VRAM).
    \item \textbf{Source Code Availability:} Full codebase publicly available for peer replication at: \url{https://github.com/mysticai11/TYPE2-PHENOTYPE}.
\end{itemize}

\section{Ethical Considerations}

\begin{itemize}[leftmargin=*]
    \item \textbf{Public De-identified Data:} The cohort was constructed entirely from public, de-identified datasets provided by the CDC. No human subject contact was performed.
    \item \textbf{Exploratory Research Prototype Only:} The system is not cleared by the FDA or CE for clinical diagnosis, and must not be used to direct patient care.
    \item \textbf{Algorithmic Bias Mitigation:} Ancestral equity analysis showed that universal clinical guidelines risk under-estimating insulin resistance in Non-Hispanic Asian populations. By conditioning the iVAE prior on demographics, LMSIS adjusts risk coordinates to demographic baselines.
    \item \textbf{No Clinical Deployment:} The FastAPI services and React dashboard are designed exclusively for demonstration and technical audit, licensed under the MIT Licence for research purposes only.
\end{itemize}

%% ============================================================
%%  THREATS TO VALIDITY
%% ============================================================
\section*{Threats to Validity}
\addcontentsline{toc}{section}{Threats to Validity}

The following table consolidates the principal threats to the validity of this study, their potential impact, and the mitigations applied.

\begin{table}[H]
\centering
\caption{Threats to validity.}
\label{tab:threats}
\begin{tabular}{p{4cm} p{4cm} p{5.5cm}}
\toprule
\textbf{Threat} & \textbf{Impact} & \textbf{Mitigation} \\ \midrule
Single survey source (NHANES only) & External validity & Temporal OOD evaluation across J- and P-cycles; zero-shot Asian cohort test \\
Cross-sectional data & No causal ordering established & All findings framed as associative; longitudinal extension deferred to future work \\
Missing CAP labels (3.7\% unlabelled; $n = 55$) & Modest loss of supervision & Semi-supervised loss masking; model successfully trains on complete cohort and handles unlabelled subset safely \\
Small Asian ancestry subgroup ($n = 12$ in reference band) & Threshold instability & Threshold marked pilot-grade; full Asian cohort ($n = 355$) confirms general shift direction \\
Survey weighting uncertainty & Wide prevalence CI & Full 95\% CI [0.00M, 64.36M] reported; estimate described as exploratory \\
Simulated pharmacological shifts & Inflated effect sizes ($r = 1.000$) & Explicitly labelled as structural disentanglement validation, not clinical trial evidence \\
Single training run (no ensemble) & Variance in results & Fixed random seed (42); integration test suite validates key properties on each run \\
\bottomrule
\end{tabular}
\end{table}

%% ============================================================
%%  CHAPTER 5: CONCLUSION AND FUTURE WORK
%% ============================================================
\chapter{Conclusion and Future Work}

\section{Conclusion}

This dissertation presented the Latent Metabolic State Inference System (LMSIS), a clinical machine-learning pipeline for detecting concurrent metabolic dysfunction in normal-BMI adults from routine blood biomarkers. The work makes five principal contributions:

\begin{enumerate}[leftmargin=*]
    \item \textbf{A Deep Generative Model Designed for Approximate Identifiability.} The \daivae{} combines demographic-conditioned priors \citep{khemakhem2020} with dual monotone anchor networks to align with theoretical iVAE identifiability conditions, with the aim of recovering approximately identifiable latent coordinates. Unlike standard VAEs, the latent axes are regularised toward insulin resistance and hepatic steatosis respectively.

    \item \textbf{Empirical Subgroup Coverage Restoration.} The implementation of Mondrian conformal prediction restores $\geq 90\%$ empirical coverage for the high-risk Dual-Burden subgroup, which standard marginal calibration leaves under-protected at 81.6\%. To the best of the authors' knowledge, following a structured literature review covering conformal prediction in healthcare, MASLD screening, metabolic phenotyping, and identifiable latent variable models from 2019--2026, this represents the first application of Mondrian conformal prediction to metabolic phenotyping.

    \item \textbf{Interpretability Extraction via Symbolic Regression.} PySR extracted closed-form approximations from the neural decoder, generating testable hypotheses about biomarker-latent relationships. The product structure in the AIP approximation suggests a hypothesis of super-additive cardiovascular risk in the Dual-Burden quadrant.

    \item \textbf{Computational Geometry on Learned Manifolds.} A Riemannian geodesic solver was implemented to plan optimal intervention paths on the latent manifold. This is presented as a computational method, not a validated clinical therapy planner.

    \item \textbf{Full-Stack Research Prototype.} The complete system---PyTorch model, FastAPI backend, React frontend, integration test suite, and public documentation---demonstrates software engineering maturity and reproducibility uncommon in student research projects.
\end{enumerate}

The system achieves a mean Spearman correlation of $\rho = 0.583 \pm 0.027$ (five-seed strict temporal mean) against CAP on the temporal OOD cohort and a combined-cohort AUROC of 0.841 against FibroScan CAP thresholds (CAP $\geq 248$\,dB/m), exceeding the performance of the evaluated clinical indices in the normal-BMI cohort. The HSI and NAFLD-LFS indices were shown to degrade structurally, with NAFLD-LFS producing risk ranking inversion in normal-BMI adults.

All results are reported with full uncertainty ranges, sample-size limitations, and simulation disclaimers. The estimated national Dual-Burden prevalence of 29.89\% ($\approx 23.91$ million) carries a correctly computed 95\% confidence interval of [0.00M, 64.36M]. This interval must accompany any citation of the point estimate.

\section{Future Work}

\begin{enumerate}[leftmargin=*]
    \item \textbf{Independent External Validation on a Real Non-NHANES Cohort.} Independent validation on a separately recruited clinical cohort, with different demographics, laboratory assays, and data-collection protocols, is the most critical gap in the current evidence base.
    \item \textbf{Prospective Outcome Validation.} Whether Dual-Burden quadrant assignment predicts longitudinal outcomes---MASLD progression, type 2 diabetes incidence, or major adverse cardiovascular events---requires prospective cohort studies.
    \item \textbf{Expanded Ancestral Equity Cohorts.} The pilot-grade finding of differential HOMA-IR thresholds across ancestral groups ($n = 12$ for Non-Hispanic Asian) requires validation in a substantially larger cohort with dedicated Asian-ancestry enrichment.
    \item \textbf{Missing-Data Robustness.} Future work should implement and evaluate missing-data robustness via partial-input inference or multiple imputation.
    \item \textbf{Longitudinal Extension.} Extending the model to longitudinal trajectories would require recurrent or state-space architectures and longitudinal training data not currently included.
    \item \textbf{Deep Learning and Tabular Baselines.} Future work should compare the system against deep tabular networks (e.g., TabNet, FT-Transformer) and other identifiable ICA architectures.
    \item \textbf{Clinician Usability Evaluation.} Formal usability studies with primary care clinicians are needed to evaluate interpretability, trust, and clinical utility of the latent-space visualisations.
    \item \textbf{Regulatory and Safety Audit.} Before any clinical deployment, the system would require FDA or CE regulatory review, adversarial robustness testing, fairness auditing, and EHR integration.
\end{enumerate}

%% ============================================================
%%  APPENDIX A: SETUP AND REPLICATION GUIDE
%% ============================================================
\appendix
\chapter{Setup and Replication Guide}

\section{Environment Setup}

Clone the repository and install the strictly version-locked core dependencies:
\begin{lstlisting}[language=bash]
git clone https://github.com/mysticai11/TYPE2-PHENOTYPE.git
cd TYPE2-PHENOTYPE
pip install -r requirements.txt
\end{lstlisting}

\section{Verify System Integrity}

Run the comprehensive integration test suite to verify model monotonicity, weight constraints, and conformal coverage bounds (15 tests):
\begin{lstlisting}[language=bash]
python -m pytest test_integration.py -v
\end{lstlisting}

\section{Launch Backend and Dashboard}

Open two separate terminal sessions:

\begin{lstlisting}[language=bash,caption={Terminal 1 - FastAPI Backend}]
python -m uvicorn main:app --app-dir backend
# API documentation: http://127.0.0.1:8000/docs
\end{lstlisting}

\begin{lstlisting}[language=bash,caption={Terminal 2 - React Dashboard}]
cd frontend
npm install
npm run dev
# Dashboard: http://localhost:5173
\end{lstlisting}

%% ============================================================
%%  REFERENCES
%% ============================================================
\bibliographystyle{apalike}
\begin{thebibliography}{99}

\bibitem[Barber et al.(2023)]{barber2023}
Barber, R.\ F., Cand\`es, E.\ J., Ramdas, A., \& Tibshirani, R.\ J. (2023).
\newblock Conformal prediction beyond exchangeability.
\newblock \textit{Annals of Statistics}, 51(2), 816--845.

\bibitem[CDC(2020)]{cdc2020}
Centers for Disease Control and Prevention. (2020).
\newblock \textit{National Health and Nutrition Examination Survey (NHANES) 2017--2020}.
\newblock U.S. Department of Health and Human Services.

\bibitem[Khemakhem et al.(2020)]{khemakhem2020}
Khemakhem, I., Kingma, D.\ P., Monti, R.\ P., \& Hyv{\"a}rinen, A. (2020).
\newblock Variational autoencoders and nonlinear ICA: A unifying framework.
\newblock \textit{Proceedings of AISTATS 2020, PMLR}, 108, 2207--2217.

\bibitem[Kingma \& Welling(2014)]{kingma2014}
Kingma, D.\ P., \& Welling, M. (2014).
\newblock Auto-encoding variational Bayes.
\newblock \textit{Proceedings of ICLR 2014}. arXiv:1312.6114.

\bibitem[Lee et al.(2024)]{lee2024}
Lee, J.\ M., et al. (2024).
\newblock Clustering-based subtyping of MASLD using liver biopsy histology.
\newblock \textit{Nature Medicine}, 30, 1--12.

\bibitem[Metwally et al.(2026)]{metwally2026}
Metwally, A.\ A., et al. (2026).
\newblock WEAR-ME: Wearable-enhanced deep learning for insulin resistance prediction in mixed-BMI cohorts.
\newblock \textit{Nature Digital Medicine}, 9, 1--14.

\bibitem[Shen et al.(2025)]{shen2025}
Shen, X., et al. (2025).
\newblock SENA-$\delta$ VAE: Causal discrepancy modelling for gene-perturbation response prediction.
\newblock \textit{Cell Systems}, 16(4), 312--328.

\bibitem[Tertulino et al.(2025)]{tertulino2025}
Tertulino, R., et al. (2025).
\newblock Random forest and XGBoost classification of MASLD from NHANES biomarkers.
\newblock \textit{PLOS ONE}, 20(3), e0287451.

\bibitem[Vovk et al.(2005)]{vovk2005}
Vovk, V., Gammerman, A., \& Shafer, G. (2005).
\newblock \textit{Algorithmic Learning in a Random World}.
\newblock Springer.

\bibitem[Zhang et al.(2025)]{zhang2025}
Zhang, Y., et al. (2025).
\newblock Machine learning prediction of MASLD in the US general population using NHANES data.
\newblock \textit{PLOS ONE}, 20(5), e0291012.

\end{thebibliography}

\end{document}
